\chapter{Introdução}
\label{chp:introduction}

A National Basketball Association (NBA) não é apenas a principal liga de basquete do mundo, mas também um negócio bilionário. Na temporada 2022-2023, o valor médio de uma franquia atingiu US\$ 3,85 bilhões, e a receita total da liga ultrapassou os US\$ 10 bilhões \citep{forbes2023}. Esse crescimento financeiro, no entanto, não acontece em um ambiente de livre mercado. A NBA opera sob um teto salarial (\textit{salary cap}) e uma \textit{luxury tax} progressiva, mecanismos criados para equilibrar competitividade e sustentabilidade comercial, embora seu efeito real ainda seja controverso, especialmente quanto à disparidade entre grandes e pequenos mercados. Na prática, grandes times conseguem montar elencos caros e muitas vezes dominantes, enquanto franquias menores precisam recorrer a estratégias mais criativas, como o desenvolvimento de jovens talentos ou a acumulação de escolhas no \textit{draft}. Essa tensão entre equilíbrio competitivo e vantagem financeira está no centro da gestão moderna da liga.

Paralelamente ao sucesso econômico, a NBA se consolidou como um dos ambientes mais avançados para o uso de dados e análises estatísticas. Sistemas de rastreamento óptico, como o SportVU e Second Spectrum, e as APIs públicas são exemplos de ferramentas que disponibilizam informações detalhadas de cada jogo, desde a posição dos jogadores em quadra até o resultado de cada arremesso \citep{oliver2004, goldsberry2019}. Esses dados permitem que times e analistas avaliem desempenho de formas que antes eram impossíveis. No entanto, transformar esse volume de informação em decisões estratégicas acertadas ainda é um desafio. Muitas organizações esbarram na dificuldade de integrar fontes distintas, limpar registros inconsistentes e organizar os dados para análises rápidas e confiáveis. Uma solução que vem ganhando espaço é a arquitetura \textit{Medallion Lakehouse}, que estrutura o processamento em camadas. Essa abordagem garante rastreabilidade e precisão, algo essencial quando se quer avaliar fenômenos complexos como o impacto de trocas de jogadores.

As trocas de elenco, as chamadas \textit{trades}, são uma das ferramentas mais utilizadas, e também as mais arriscadas, da gestão esportiva. Diferente da contratação de agentes livres na entressafra, uma troca geralmente ocorre com a temporada em andamento. O novo jogador precisa se adaptar imediatamente a um novo sistema tático, a novos companheiros e, muitas vezes, a uma nova cidade. A frequência dessas movimentações cresceu sensivelmente nas últimas temporadas. Considerando apenas a janela de 30 dias que antecede o \textit{trade deadline}, recorte temporal adotado nesta pesquisa, o número de jogadores movimentados passou de 35 na temporada 2012-13 para 107 na temporada 2024-25, um crescimento superior a 200\% em treze anos (Figura~\ref{fig:volume_trocas}). Esse cenário torna urgente compreender de forma rigorosa se e como as trocas afetam o desempenho das equipes.

\begin{figure}[htb]
    \centering
    \caption{Volume de jogadores trocados na janela de 30 dias antes do \textit{trade deadline} (2012-13 a 2024-25)}
    \label{fig:volume_trocas}
    \includegraphics[width=.85\textwidth]{images/introducao_volume_trocas_temporada.png}
    \begin{flushleft}
    \footnotesize Fonte: elaborada pela autora (2026).
    \end{flushleft}
\end{figure}

Diante disso, o problema central que motiva esta pesquisa é justamente a dificuldade de antecipar se uma troca trará ganho real de performance para a equipe receptora. As decisões de troca costumam se basear em métricas de produtividade individual, como o \textit{Player Efficiency Rating} (PER) desenvolvido por \citet{hollinger2007}, no entanto, esse tipo de medida ignora ao menos duas dimensões importantes do valor de um jogador para o time que o adquire. A primeira é a adequação funcional, pois pode acontecer do novo jogador suprir uma deficiência específica do time, como rebotes, eficiência defensiva ou aproveitamento de três pontos, ou ele apenas adicionar talento a uma posição já bem servida. A segunda é o impacto da troca sobre a distribuição salarial interna, essa que altera a folha de pagamentos, e que pode afetar a percepção de justiça no vestiário e, consequentemente, a coesão e a motivação do grupo. Além dessas duas lacunas na avaliação do jogador em si, a própria estrutura da negociação carrega informação raramente considerada: jogadores mais experientes tendem a precisar de menos tempo de adaptação ao novo sistema, e lacunas posicionais mais severas tendem a exigir pacotes de troca maiores para serem fechadas. Cada uma dessas dimensões exige uma lógica de avaliação diferente, e ignorá-las empobrece a leitura do que, de fato, torna uma troca bem-sucedida.

A literatura econômica apresenta duas perspectivas concorrentes sobre o efeito da desigualdade salarial no desempenho coletivo. Uma delas é a Teoria da Equidade \citep{adams1965} que sustenta que disparidades excessivas geram percepções de injustiça, reduzem a cooperação e prejudicam os resultados. Já a Teoria do Torneio \citep{lazear1981} argumenta que grandes diferenças salariais funcionam como incentivo à produtividade, refletindo hierarquias de talento e estimulando o esforço individual. Os estudos empíricos na NBA não chegam a um consenso. Enquanto alguns autores não encontraram efeito significativo da dispersão salarial \citep{berri2004, katayama2011}, outros identificaram impactos positivos \citep{frick2003, simmons2011}. Mais recentemente, \citet{ching2019} sugeriu uma relação em forma de U invertido, onde até certo ponto, a desigualdade estimula o desempenho; além desse limite, os efeitos negativos sobre a coesão se sobrepõem. Além disso, a mensuração da desigualdade pode ser aprimorada pelo uso do coeficiente de Gini ajustado, que considera o limite superior factível de vitórias em uma liga esportiva \citep{larsen2006}.

Outra tensão relevante diz respeito à própria lógica da troca. A Teoria do \textit{O-ring} \citep{kremer1993}, aplicada ao basquetebol por \citet{ching2019}, propõe que a produção em equipe depende do elo mais fraco. Em outras palavras, times com talento homogeneamente distribuído podem superar aqueles que concentram recursos em uma única superestrela e têm um elenco deficiente. No basquetebol, a complementaridade funcional, isto é, adicionar um jogador que supre uma deficiência específica, tende a ter impacto mais significativo do que a simples adição de uma estrela em posição já bem servida \citep{oliver2004, kubatko2007}. Uma forma de medir essa complementaridade na prática é o \textit{Win Probability Added} (WPA), que quantifica o impacto de um jogador pela variação que suas ações produzem na probabilidade de vitória do time ao longo da partida, originalmente proposto para jogadores da NBA por \citet{deshpande2016}; é sobre essa lógica que se apoia a Contribuição Relativa \citep{sa2025}, metodologia adotada nesta pesquisa para diagnosticar, por grupo posicional, se uma troca de fato preenche uma lacuna estrutural do elenco ou apenas adiciona talento a um grupo posicional cuja contribuição já está na média da liga ou acima dela.

O contexto competitivo da equipe também pode moderar os efeitos das trocas. \citet{katayama2011} mostram que o impacto da dispersão salarial não é uniforme entre times com diferentes probabilidades de chegar aos \textit{playoffs}, e os chamados \textit{tournament incentives} \citep{taylor2002} sugerem que times fortes e fracos podem ter incentivos opostos dentro de uma mesma temporada. Esse fator raramente é incorporado nos modelos de impacto de trocas, mas ajuda a explicar por que um mesmo efeito pode se concentrar em alguns perfis de time e não em outros.

Diante desse cenário de controvérsias empíricas e lacunas teóricas, como a ausência de consenso sobre o efeito da dispersão salarial, pouca atenção ao ajuste funcional dos jogadores, rara consideração do perfil competitivo das franquias e limitações metodológicas das análises agregadas por temporada, torna-se evidente que o impacto das trocas de elenco na NBA ainda não foi devidamente compreendido. As decisões de \textit{front office} continuam sendo tomadas, muitas vezes, com base em intuição ou em métricas individuais que ignoram a complexidade do fenômeno. Esta pesquisa parte da premissa de que é preciso tratar as trocas como um evento multidimensional, no qual aspectos financeiros, táticos e estratégicos se entrelaçam. Ao fazer isso, espera-se oferecer uma contribuição empírica que ajude a reduzir as incertezas que cercam uma das ferramentas mais arriscadas, e ao mesmo tempo mais utilizadas, da gestão esportiva moderna.

\section{Objetivos}
\label{sec:objectives}
O desenvolvimento deste trabalho tem como foco a construção de um modelo multidimensional para analisar o impacto das trocas de elenco na performance de times da NBA. Para isso, utiliza-se uma arquitetura de dados em nível de jogo, métodos estatísticos de correlação e ferramentas de visualização gráfica para identificar padrões e interpretar os resultados. A Figura~\ref{fig:modelo_multidimensional} resume as quatro dimensões que compõem esse modelo, cada uma retomada nos objetivos específicos a seguir.

\begin{figure}[htb]
    \centering
    \caption{Modelo multidimensional de análise do impacto das trocas}
    \label{fig:modelo_multidimensional}
    \includegraphics[width=.85\textwidth]{images/introducao_modelo_multidimensional.png}
    \begin{flushleft}
    \footnotesize Fonte: elaborada pela autora (2026).
    \end{flushleft}
\end{figure}

\subsection{Objetivo Geral}
Analisar o impacto das trocas de elenco na performance de times da NBA por meio de um modelo multidimensional que considere a disparidade salarial em quadra, o ajuste funcional do jogador adquirido e as características da negociação, como a experiência dos jogadores envolvidos e o tamanho do pacote de troca.

\subsection{Objetivos Específicos}
\label{sec:specific_objectives}
Para viabilizar a construção desse modelo e garantir a qualidade das análises, foram definidos os seguintes objetivos específicos:
\begin{itemize}
    \item Coletar dados estatísticos e contratuais da NBA por meio de APIs públicas e fontes históricas, implementando a arquitetura \textit{Medallion Lakehouse} para organizar os dados nas camadas Bronze, Silver e Gold;
    \item Projetar a modelagem dimensional dos dados e orquestrar o \textit{pipeline} de processamento utilizando SQL para consultas e transformações estruturais, e Python para automação, integração entre etapas e garantia de reprodutibilidade e rastreabilidade analítica;
    \item Avaliar seis métricas candidatas de disparidade salarial: desvio padrão, coeficiente de variação, razão entre o jogador mais bem pago e o menos bem pago do elenco, proporção da folha salarial concentrada nos três maiores contratos, Índice de Gini simples e Índice de Gini ponderado por minutos jogados, adotando esta última como métrica principal por ser a única que pondera cada jogador pelos minutos efetivamente jogados, refletindo a desigualdade que o elenco vive em quadra e não apenas a distribuição nominal da folha salarial, e confirmando essa escolha por meio de correlação de Spearman e diagramas de dispersão com a variação de desempenho após a troca;
    \item Verificar se essa associação entre desigualdade salarial e desempenho é homogênea entre os arquétipos competitivos do time, obtidos por clusterização (\textit{k-means}), de modo a identificar em quais perfis de equipe o efeito observado se concentra;
    \item Avaliar o impacto técnico das trocas a partir da Contribuição Relativa, métrica derivada do \textit{Win Probability Added} (WPA) que mede, por grupo posicional, se a aquisição preenche uma lacuna estrutural do elenco ou apenas reforça uma posição já bem servida \citep{sa2025}, e da variação do \textit{Net Rating} ajustado pela força do adversário em janelas temporais pré e pós-evento, empregando \textit{boxplots} e gráficos de evolução pré e pós-troca para identificar padrões de melhora ou declínio;
    \item Testar se a experiência dos jogadores adquiridos, medida pela idade média no momento da troca, está associada ao impacto de desempenho de curto prazo do time, por meio de correlação de Spearman e comparação por quartis de idade;
    \item Testar se a severidade da lacuna posicional preenchida está associada ao tamanho do pacote de jogadores envolvidos na negociação, por meio de correlação de Spearman e comparação por quartis de severidade.
\end{itemize}

A partir desses objetivos, busca-se desenvolver uma análise robusta e empiricamente fundamentada, capaz de oferecer evidências tanto para o debate acadêmico sobre teorias econômicas aplicadas ao esporte quanto para a prática de gestores e analistas da NBA.

\section{Justificativa}
A literatura especializada ainda apresenta resultados divergentes sobre os efeitos das trocas de elenco no desempenho das equipes da NBA. Enquanto estudos como o de \citet{larsen2006} demonstram que a mensuração da desigualdade salarial pode ser aprimorada pelo uso do coeficiente de Gini ajustado, com implicações diretas para a análise do impacto de trocas, outros trabalhos encontram efeitos positivos, negativos ou mesmo inexistentes, evidenciando a falta de consenso na área. Essa controvérsia não é meramente acadêmica, ela reflete a dificuldade real que gestores de franquias enfrentam ao decidir se e quando realizar uma troca, e sob quais condições ela tende a ser bem-sucedida.

Grande parte dos estudos existentes concentra-se em análises agregadas por temporada, tratando a dispersão salarial como uma característica estática do elenco. No entanto, como apontam \citet{katayama2011}, a estrutura salarial de uma equipe pode variar significativamente ao longo da temporada, especialmente devido a trocas de meio de ano. Além disso, a maioria dos modelos não incorpora o ajuste funcional, isto é, se o jogador adquirido supre uma deficiência específica do time, auxiliando nos rebotes, arremessos de três pontos ou eficiência defensiva, ou se ele apenas adiciona talento a uma posição já bem servida. Dessa forma, ignorar o ajuste funcional é especialmente problemático porque, em um ambiente de teto salarial apertado, times não podem desperdiçar contratações com talentos redundantes.

Outra lacuna relevante diz respeito à heterogeneidade entre os times. Como mostram \citet{taylor2002}, os incentivos para vencer ou perder não são uniformes entre franquias, e poucos modelos de impacto de trocas verificam se um mesmo efeito se sustenta igualmente entre diferentes perfis competitivos. O que nos leva a entender que tratar a liga como um bloco homogêneo pode levar a conclusões enganosas, já que um mesmo tipo de troca pode ter efeitos distintos dependendo do estilo de jogo e do momento competitivo do time que a realiza.

Diante desse panorama, este trabalho se justifica por abordar as trocas de elenco sob uma perspectiva multidimensional, incorporando variáveis como o ajuste funcional dos jogadores às lacunas estatísticas do time, a desigualdade salarial medida pelo índice de Gini, a experiência dos jogadores envolvidos e a estrutura da negociação. Metodologicamente, a proposta se diferencia ao utilizar uma arquitetura de dados em nível de jogo, por meio de um \textit{pipeline} estruturado em camadas, uma para os dados brutos, outra para a limpeza e validação, e a última para a modelagem dimensional. Essa abordagem permite capturar variações intra temporada, como a participação efetiva de cada jogador em cada partida, que análises agregadas por temporada não conseguem detectar. Além disso, ao empregar técnicas de aprendizado de máquina, como \textit{k-means} e análise de componentes principais, e janelas temporais móveis (pré e pós-troca), é possível isolar com mais precisão o efeito da movimentação, ajustando o desempenho pela força do adversário enfrentado e testando a robustez dos resultados a diferentes larguras de janela ao redor do \textit{trade deadline}.

Como resultado, espera-se que esta pesquisa ofereça de forma integrada as previsões concorrentes da Teoria da Equidade e da Teoria do Torneio no contexto específico das trocas na NBA, algo que estudos anteriores fizeram apenas de maneira fragmentada. Que ofereça também uma arquitetura de dados em nível de jogo que possa ser aplicada a um problema de gestão esportiva, gerando evidências mais finas do que as análises sazonais tradicionais. E por fim, que os resultados auxiliem os \textit{front offices} a avaliar trocas com base não apenas no talento bruto, mas também na adequação funcional, na experiência dos jogadores envolvidos e na estrutura da negociação, reduzindo a incerteza em decisões que envolvem milhões de dólares.

\section{Estrutura do Trabalho}

Este trabalho está estruturado em cinco capítulos. Depois desta introdução, a fundamentação teórica é apresentada no Capítulo 2, que aborda os fundamentos do \textit{sports analytics}, as métricas avançadas de basquetebol, as teorias de gestão de elencos, os estudos empíricos anteriores sobre dispersão salarial, trocas e desempenho na NBA, além da arquitetura de dados \textit{Medallion Lakehouse} que sustenta a coleta e o tratamento das informações utilizadas. Na sequência, o Capítulo 3 volta-se à metodologia, detalhando a construção da amostra e a operacionalização de cada hipótese de pesquisa em variáveis mensuráveis e procedimentos estatísticos, sem deixar de lado as variáveis moderadoras e de contexto e as limitações identificadas. Os resultados obtidos vêm no Capítulo 4, relacionados às hipóteses de pesquisa e à literatura discutida anteriormente. Por fim, o Capítulo 5 encerra o trabalho com as considerações finais, as limitações do estudo, as implicações gerenciais e as sugestões para pesquisas futuras.
